\emph{(i)} $P_{12}$ fixes $u'$ and $\hat{b}_{e}$ and reverses $\hat{b}_{o}$, so for $P_{12}XP_{12} = -X$ every entry $\hat{b}^{\mathsf{T}}X\hat{b}'$ with $\hat{b}, \hat{b}'$ both in $\{u', \hat{b}_{e}\}$ or both equal to $\hat{b}_{o}$ is equal to its own negative and vanishes. For $E$: $\hat{b}_{e}^{\mathsf{T}}E\hat{b}_{e} = (e_{3} \cdot \hat{b}_{e})^{2} = \tfrac23$ and $E\hat{b}_{o} = 0$. Hence the compression of $E + X_{s}$ to $u'^{\perp}$ is the displayed block, with $x_{s} = \hat{b}_{e}^{\mathsf{T}}X_{s}\hat{b}_{o}$ and $y_{s} = \hat{b}_{o}^{\mathsf{T}}X_{s}\hat{b}_{e}$, which are equal for symmetric and opposite for antisymmetric $X_{s}$. Since $\Pi^{2} = \Pi$, $\Pi A\Pi A^{\mathsf{T}}\Pi = (\Pi A\Pi)(\Pi A^{\mathsf{T}}\Pi) = (\Pi A\Pi)(\Pi A\Pi)^{\mathsf{T}}$. By Lemma \claimref{neutrino.breaking-relation}, at its order the second and third light neutrino masses are $(d^{2}\delta^{2}/\alpha)$ times the eigenvalues of this compression, with its eigenvectors as the light states; they are the squared singular values and the left singular vectors of $B_{\nu}$. For the charged leptons, $M_{e} = c\,u'u'^{\mathsf{T}} + c\delta_{c}(E + X_{e})$ is a rank-one matrix plus a perturbation whose left and right kernels are $u'^{\perp}$, and at the same order its two light singular values and left singular vectors are $c\delta_{c}$ times those of the compression to $u'^{\perp} \times u'^{\perp}$, which is $B_{e}$; this is the charged-lepton reading the Lemma's instances use (Proposition \claimref{synthesis.theta23-at-s3-s2-direction}).

\emph{(ii)} Write the singular value decomposition $B = R(\theta_{L})\,\mathrm{diag}(\sigma_{1}, s_{2})\,R(\theta_{R})^{\mathsf{T}}$ with $R$ the rotation matrix and $|s_{2}| = \sigma_{2}$ (a reflection is absorbed in the sign of $s_{2}$). Its $(2,2)$ entry is $\sigma_{1}\sin\theta_{L}\sin\theta_{R} + s_{2}\cos\theta_{L}\cos\theta_{R} = 0$, so $s_{2}/\sigma_{1} = -\tan\theta_{L}\tan\theta_{R}$ (no cosine vanishes, since $\cos\theta_{L}\cos\theta_{R} = 0$ would force $\sigma_{1}\sin\theta_{L}\sin\theta_{R} = 0$ and make $B$ rank one or $k = 0$). The heavy left and right singular vectors are the heavy eigenvectors of $BB^{\mathsf{T}} = \bigl[\begin{smallmatrix} k^{2} + x^{2} & ky \\ ky & y^{2} \end{smallmatrix}\bigr]$ and $B^{\mathsf{T}}B = \bigl[\begin{smallmatrix} k^{2} + y^{2} & kx \\ kx & x^{2} \end{smallmatrix}\bigr]$, whose tangents solve $ky\,t^{2} + (k^{2} + x^{2} - y^{2})t - ky = 0$ and $kx\,t^{2} + (k^{2} + y^{2} - x^{2})t - kx = 0$. For $y = x$ the two quadratics coincide and for $y = -x$ they are exchanged by $t \mapsto -t$; each has the two roots $t$ and $-1/t$, the heavy and the light direction, and the heavy one is the same choice in both since the singular values are shared. So $\theta_{L} = \pm\theta_{R}$.

\emph{(iii)} In the light plane the lightest state of each sector is orthogonal to its heavy light-plane state, so the angle between the lightest states equals the angle between the heavy ones, $|\theta_{L}^{\nu} - \theta_{L}^{e}|$, and $|U_{e1}| = |\langle e, \nu_{1}\rangle| = \cos\Delta$. For one-type second stages, (ii) gives $\tan^{2}\theta_{L}^{e} = (\sigma_{2}/\sigma_{1})_{e} = m_{e}/m_{\mu}$ and $\tan^{2}\theta_{L}^{\nu} = (\sigma_{2}/\sigma_{1})_{\nu} = (m_{1}/m_{2})^{1/2}$, whence the bound, by $|\theta_{L}^{\nu} - \theta_{L}^{e}| \le |\theta_{L}^{\nu}| + |\theta_{L}^{e}|$. Numerically, with $a = 0.06955 \ge \sqrt{m_{e}/m_{\mu}}$ at the PDG masses and $b = 0.31623 \ge 0.01^{1/4}$, $\tan\Delta \le (a + b)/(1 - ab) < 0.3959 < \tan 21.6^{\circ}$ and $\cos^{2}\Delta > 0.865$. The global fit's three-sigma edges $s_{12}^{2} \ge 0.275$, $s_{13}^{2} \ge 0.0203$ give $|U_{e1}|^{2} \le 0.725 \times 0.9797 < 0.711$.

\emph{(iv)} $\cos^{2}\Delta \le 0.711$ gives $\tan\Delta \ge 0.6375$. If the charged sector is of one type, $|\theta_{L}^{e}| \le \arctan a$, so $\tan|\theta_{L}^{\nu}| \ge (0.6375 - a)/(1 + 0.6375a)$, and by (ii) $|\tan\theta_{L}^{\nu}\tan\theta_{R}^{\nu}| = (m_{1}/m_{2})^{1/2} \le \tfrac{1}{10}$ when $m_{1}/m_{2} \le 0.01$, so $|\tan\theta_{L}^{\nu}/\tan\theta_{R}^{\nu}| \ge 10\tan^{2}\theta_{L}^{\nu} > 2.95$. For the texture entries, the decomposition of (ii) gives $x = \sigma_{1}\cos\theta_{L}\sin\theta_{R} - s_{2}\sin\theta_{L}\cos\theta_{R}$ and $y = \sigma_{1}\sin\theta_{L}\cos\theta_{R} - s_{2}\cos\theta_{L}\sin\theta_{R}$; with $s_{2}/\sigma_{1} = -\tan\theta_{L}\tan\theta_{R}$ and dividing by $\sigma_{1}\cos\theta_{L}\cos\theta_{R}$, exactly $y/x = t_{L}(1 + t_{R}^{2})/\bigl(t_{R}(1 + t_{L}^{2})\bigr)$ with $t = \tan\theta$. Writing $\rho = |t_{L}t_{R}| = (m_{1}/m_{2})^{1/2}$, $|y/x| = (t_{L}^{2} + \rho^{2})/\bigl(\rho(1 + t_{L}^{2})\bigr)$, which increases with $t_{L}^{2}$ (its derivative has the sign of $1 - \rho^{2}$) and, for $\rho^{2} < t_{L}^{2}$, decreases with $\rho$; at $t_{L}^{2} > 0.295$ and $\rho \le \tfrac{1}{10}$ it exceeds $0.305/0.1295 > 2.35$. (The first-order relation $|\tan\theta_{L}/\tan\theta_{R}| \approx |y/x|$, from $\tan 2\theta_{L} = 2ky/(k^{2} + x^{2} - y^{2})$ and $\tan 2\theta_{R} = 2kx/(k^{2} + y^{2} - x^{2})$, is only a small-$x$, small-$y$ guide and is not used.) Since $|y/x| > 1$, $|y| \ne |x|$, which requires both a symmetric part $(x + y)/2$ and an antisymmetric part $(y - x)/2$.

The script verifies (i) on rational $S_{2}$-odd, symmetric and antisymmetric second stages, (ii) on singular value decompositions built from rational points of the circle (the zero corner, the heavy vectors, the tilt quadratics, the ratio identity and the one-type coincidence), the exact relation of $y/x$ to the tilts on those decompositions, and the numerical bounds of (iii) and (iv) as exact rational inequalities.
